\section{为什么需要向量空间}
\label{sec:03-Linear-Algebra/02-Vector-Spaces/01}

\subsection{一个到处重演的模式}

你正在解一个齐次线性系统：

\[
\begin{pmatrix}
1 & 2 & -1 \\
2 & 4 & 1 \\
-1 & -2 & 2
\end{pmatrix}
\begin{pmatrix}x\\y\\z\end{pmatrix}
=
\begin{pmatrix}0\\0\\0\end{pmatrix}.
\]

消元之后你找到了两个不共线的解：

\[
z_1=\begin{pmatrix}-2\\1\\0\end{pmatrix},\qquad
z_2=\begin{pmatrix}1\\0\\1\end{pmatrix}.
\]

然后你注意到一件有意思的事：\(z_1+z_2=(-1,1,1)^{\trans}\)，代回去，也是解。\(3z_1=(-6,3,0)^{\trans}\)，还是解。\(2z_1-5z_2=(-9,2,-5)^{\trans}\)，依然是解。

这不是碰巧。如果你把两个解各自代入 \(Ax=0\)，然后相加：

\[
A(z_1+z_2)=Az_1+Az_2=0+0=0.
\]

如果你把一个解放大 \(c\) 倍：

\[
A(cz_1)=cAz_1=c\cdot0=0.
\]

解与解的任意线性组合，仍然是解。齐次线性系统的解集不是一组散乱的点——它们构成一个对加法和数乘\emph{封闭}的整体。

你可能会想：这只是线性方程组中的一个小性质而已，有必要专门拿出来讲吗？

换个场景。考虑微分方程

\[
y''+y=0.
\]

你知道 \(\sin t\) 是一个解，\(\cos t\) 也是一个解。那么 \(3\sin t-2\cos t\) 呢？代入：

\[
(3\sin t-2\cos t)''+(3\sin t-2\cos t)
=(-3\sin t+2\cos t)+(3\sin t-2\cos t)=0.
\]

也是解。任意线性组合都是解。

再换一个场景。次数不超过 2 的多项式：\(p(t)=a+bt+ct^2\)。两个这样的多项式相加，次数不会超过 2。数乘之后，次数仍然不超过 2。线性组合始终留在这个集合里。

三个不同的数学对象——\(\R^n\) 中的向量，微分方程的解，多项式——做的事情却完全一样：集合中的元素可以进行线性组合，组合结果仍留在集合中。

如果你在三个不同的地方看到同一个模式，而每次都要重新证明一遍``组合之后还在里面''，那说明这个模式值得被抽象出来，给它一个名字和一套统一的推理工具。

\subsection{向量不只是一支箭头}

在 \(\R^n\) 中，向量是 \(n\) 个有序实数。画在纸上的箭头只是它的几何形象。但只要定义了合适的加法和数乘规则，以下对象全部可以当作向量来操作：

\begin{itemize}
\tightlist
\item 次数不超过 2 的多项式 \(a+bt+ct^2\)——加法就是对应系数相加，数乘就是每个系数乘以标量；
\item 从 \([0,1]\) 到 \(\R\) 的连续函数——函数的加法逐点定义，数乘也逐点定义；
\item \(m\times n\) 矩阵——矩阵加法逐元素进行，数乘每个元素都乘；
\item 一段有限长度的离散信号 \((s_1,\ldots,s_N)\)；
\item 某个齐次线性微分方程的全部解；
\item 某个齐次线性方程组 \(Ax=0\) 的全部解。
\end{itemize}

这些对象外形千差万别，但共享一套运算规则：加法结合律（\((u+v)+w=u+(v+w)\)）、加法交换律（\(u+v=v+u\)）、存在零元素（\(0+u=u\)）、存在加法逆元（每个 \(u\) 有一个 \(-u\) 使得 \(u+(-u)=0\)）、数乘对加法的分配律（\(c(u+v)=cu+cv\) 和 \((c+d)u=cu+du\)）、数乘结合律（\(c(du)=(cd)u\)），以及 \(1\cdot u=u\)。

满足这八条规则的集合，连同定义在其上的加法和数乘，称为一个\emph{向量空间}。

八条规则看起来很啰嗦，但你不需要逐条背诵。你只需要理解它们共同保证的一件事：\emph{线性组合可以自由地重新分组、缩放、展开，并且结果始终是同一种对象。}后面所有关于基、维数、线性变换的结论，都建立在这个保证之上。公理不是凭空制定的教条——它们是让线性组合这种操作``合法化''的最低要求。

如果你好奇某条公理能不能从其他几条推导出来：加法交换律确实可以从其余的公理推导（利用分配律和 \(1+1=2\)），但在有限域上的向量空间中这个推导不成立，所以通常保留所有八条以确保普适性。这个细节属于代数课程，此处无需深究。

\subsection{子空间：大空间里面的封闭小世界}

设 \(V\) 是一个向量空间。如果它的一个非空子集 \(S\subseteq V\) 自身也满足向量空间的全部要求，就称 \(S\) 是 \(V\) 的\emph{子空间}。

你不用逐条验证八条公理。绝大多数公理（结合律、交换律、分配律）直接从 \(V\) 继承——只要运算规则一样，它们在子集里自然成立。你需要检查的只有一件事：\(S\) 是否对加法和数乘\emph{封闭}。

确切地说，下面这一个条件就够了：

\[
\text{对任意 } u,v\in S \text{ 和任意标量 } \alpha,\beta,\quad
\alpha u+\beta v\in S.
\]

为什么这一个条件就包含了所有需要？取 \(\alpha=\beta=0\)，得到 \(0\in S\)（零向量必须在里面）。取 \(\alpha=1,\beta=0\)，得到加法封闭。取 \(\beta=0\)，得到数乘封闭。加法逆元可以从数乘得到：取 \(\alpha=-1\)，则 \(-u=(-1)u\in S\)。

用一个具体的子空间来感受一下。在 \(\R^3\) 中，考虑所有满足 \(x+y+z=0\) 的点：

\[
S=\left\{
\begin{pmatrix}x\\y\\z\end{pmatrix}
:x+y+z=0
\right\}.
\]

这个集合是 \(\R^3\) 中穿过原点的平面。检验封闭性：设 \(u=(u_1,u_2,u_3)^{\trans}\) 和 \(v=(v_1,v_2,v_3)^{\trans}\) 都在 \(S\) 中，即它们的坐标和都为零。那么 \(\alpha u+\beta v\) 的坐标和为

\[
(\alpha u_1+\beta v_1)+(\alpha u_2+\beta v_2)+(\alpha u_3+\beta v_3)
=\alpha(u_1+u_2+u_3)+\beta(v_1+v_2+v_3)=0.
\]

通过验证。这是一个子空间。

再看一个反例：

\[
T=\left\{
\begin{pmatrix}x\\y\\z\end{pmatrix}
:x+y+z=1
\right\}.
\]

这也是一个平面，但不穿过原点。零向量 \((0,0,0)^{\trans}\) 不满足 \(x+y+z=1\)，因而 \(0\notin T\)，直接排除了子空间资格。实际上，取 \(u=(1,0,0)^{\trans}\in T\) 和 \(v=(0,1,0)^{\trans}\in T\)，它们的和是 \((1,1,0)^{\trans}\)，坐标和为 \(2\neq1\)，不在 \(T\) 中。加法的封闭性也被破坏。

\(T\) 不是一个向量空间，但它是子空间 \(S\) 加上一个平移向量得到的结果。这类集合称为\emph{仿射子空间}——后面会在非齐次解集中再次遇到。

\subsection{线性组合与张成：用几个方向生成一个世界}

给定向量空间中的一组向量 \(v_1,\ldots,v_k\)，取任意标量 \(c_1,\ldots,c_k\) 形成

\[
c_1v_1+\cdots+c_kv_k,
\]

得到的结果称为 \(v_1,\ldots,v_k\) 的\emph{线性组合}。所有可能的线性组合构成的集合记作

\[
\operatorname{span}\{v_1,\ldots,v_k\}.
\]

称为由 \(v_1,\ldots,v_k\)\emph{张成}（span）的空间。

张成空间有一个``最小性''：它是包含 \(v_1,\ldots,v_k\) 的最小子空间。任何包含这些向量的子空间，因为必须对线性组合封闭，一定包含它们的所有组合，因此一定包含整个 \(\operatorname{span}\)。你不能在比 \(\operatorname{span}\) 更小的子空间里装下这些向量。

回到 \(\R^3\) 中的平面 \(x+y+z=0\)。取两个向量：

\[
v_1=\begin{pmatrix}1\\-1\\0\end{pmatrix},\qquad
v_2=\begin{pmatrix}1\\0\\-1\end{pmatrix}.
\]

这两个向量的坐标和都为零，所以它们属于 \(S\)。它们的任意线性组合 \(c_1v_1+c_2v_2\) 的坐标和也是零——因为每个 \(v_i\) 的坐标和为零，线性组合的坐标和就是 \(c_1\cdot0+c_2\cdot0=0\)。所以

\[
\operatorname{span}\{v_1,v_2\}\subseteq S.
\]

反过来，任取 \((x,y,z)^{\trans}\in S\)，有 \(z=-x-y\)。可以将它表示为 \(v_1,v_2\) 的组合：

\[
\begin{pmatrix}x\\y\\-x-y\end{pmatrix}
=(-y)\begin{pmatrix}1\\-1\\0\end{pmatrix}
+(x+y)\begin{pmatrix}1\\0\\-1\end{pmatrix}.
\]

所以 \(S\subseteq\operatorname{span}\{v_1,v_2\}\)。两边互相包含，因此 \(S=\operatorname{span}\{v_1,v_2\}\)。

同一个平面，你有了两种描述方式：

\begin{itemize}
\tightlist
\item \emph{隐式描述}（方程）：向量必须满足 \(x+y+z=0\)；
\item \emph{参数描述}（生成方向）：向量可以写成 \(v_1,v_2\) 的线性组合。
\end{itemize}

隐式描述告诉你\emph{限制}是什么，参数描述告诉你\emph{自由度}在哪里。Gaussian 消元的核心功能，正是在这两种描述之间做转换：从方程出发，解出自由变量，写出参数形式。

\subsection{列空间：变换能到达的所有地方}

矩阵

\[
A=\begin{pmatrix}a_1&\cdots&a_n\end{pmatrix}
\]

的列空间定义为

\[
\operatorname{Col}(A)=\operatorname{span}\{a_1,\ldots,a_n\}.
\]

它包含了所有可能的输出 \(Ax\)（因为 \(Ax=x_1a_1+\cdots+x_na_n\) 正是列的线性组合）。

列空间天然是子空间：任意两个输出的线性组合可以通过分配律收回矩阵乘法内部：

\[
\alpha Ax+\beta Ay=A(\alpha x+\beta y),
\]

而 \(\alpha x+\beta y\) 仍是某个输入，所以 \(\alpha Ax+\beta Ay\) 也是 \(A\) 的某个输出，仍在 \(\operatorname{Col}(A)\) 中。

这使得可解性有了一个极其简洁的表达：

\[
Ax=b\text{ 有解}
\quad\Longleftrightarrow\quad
b\in\operatorname{Col}(A).
\]

你不需要解方程就知道答案是否存在——只要看 \(b\) 是否在列空间里。消元中的矛盾行 \([0\;\cdots\;0\mid c]\)，\(c\neq0\)，就是用代数语言告诉你：\(b\) 偏离了列空间，多了一个列空间够不到的分量。

\subsection{零空间：变换丢掉了哪些方向}

矩阵的零空间定义为

\[
\mathcal N(A)=\{x:Ax=0\}.
\]

它包含所有被 \(A\)``压扁''到零向量的输入。本章开头的计算已经确认了零空间对线性组合封闭：

\[
A(\alpha z_1+\beta z_2)=\alpha Az_1+\beta Az_2=0.
\]

所以零空间也是输入空间的子空间。

零空间的几何意义是：如果 \(z\in\mathcal N(A)\) 且 \(z\neq0\)，那么沿 \(z\) 方向移动输入，输出纹丝不动：

\[
A(x+z)=Ax+Az=Ax.
\]

这个方向上的信息被变换完全丢弃了。零空间非平凡（维度大于 0）意味着变换不是一对一的：多个不同的输入被映射到同一个输出。

列空间和零空间承担完全不同的角色，但它们之间存在深刻的联系。列空间位于\emph{输出}空间，回答``能到达哪里''（可达性）。零空间位于\emph{输入}空间，回答``丢掉了什么''（信息损失）。后面单元的秩-零化度定理会把两者用一个等式连起来：输入维数 = 秩（列空间维数） + 零空间维数。

\subsection{为什么非齐次解集不是子空间：平移后的世界}

设 \(S_b=\{x:Ax=b\}\)，其中 \(b\neq0\)。首先零向量不在 \(S_b\) 中，因为 \(A0=0\neq b\)。连零都没有，不可能是子空间。

再看加法封闭性。取 \(x_1,x_2\in S_b\)：

\[
A(x_1+x_2)=Ax_1+Ax_2=b+b=2b\neq b.
\]

两个解的和变成了 \(2b\) 的解，不是 \(b\) 的解。所以非齐次解集对加法不封闭。

它的结构我们在消元部分已经见过：如果 \(x_p\) 是一个特解，那么全体解为

\[
S_b=x_p+\mathcal N(A)=\{x_p+z:z\in\mathcal N(A)\}.
\]

你拿到的是一个``子空间 + 平移''的形状——在几何上它和零空间有同样的方向和维度，只是整个结构被抬离了原点。这类集合称为仿射子空间（affine subspace）。

区分齐次解集和非齐次解集不是咬文嚼字。齐次解集满足叠加原理——你可以把解线性组合得到新解。非齐次解集不满足——两个解的简单平均可能已经不是解。在实际问题中，如果你误以为可以对非齐次解做线性组合，得到的``解''代回原方程就会露馅。

\subsection{多项式中的子空间：没有箭头的向量空间}

脱离 \(\R^n\) 的几何直觉，用多项式做一个例子。令 \(P_2\) 表示所有次数不超过 2 的实多项式：

\[
P_2=\{a+bt+ct^2:a,b,c\in\R\}.
\]

加法和数乘按逐系数进行：\((a+bt+ct^2)+(a'+b't+c't^2)=(a+a')+(b+b')t+(c+c')t^2\)。这满足向量空间的所有公理——多项式的加法逐系数进行，和系数向量的加法完全同步。

现在考虑一个子集：

\[
S=\{p\in P_2:p(1)=0\}.
\]

在 \(t=1\) 处取值为零的多项式。检验子空间条件：取 \(p,q\in S\)，即 \(p(1)=q(1)=0\)，则

\[
(\alpha p+\beta q)(1)=\alpha p(1)+\beta q(1)=0.
\]

通过。\(S\) 是 \(P_2\) 的子空间。

每个满足 \(p(1)=0\) 的多项式必然含有因子 \((t-1)\)（这是多项式的一个基本性质：\(p(a)=0\) 当且仅当 \((t-a)\) 整除 \(p(t)\)）。由于 \(p\) 的次数不超过 2，它可以写成

\[
p(t)=(t-1)(r+st)=r(t-1)+s(t^2-t).
\]

也就是说，\(S\) 中的任意多项式都可以用 \(t-1\) 和 \(t^2-t\) 这两个``基多项式''线性组合出来：

\[
S=\operatorname{span}\{t-1,\,t^2-t\}.
\]

这和前面 \(\R^3\) 中的平面 \(x+y+z=0\) 是同构的结构：一个齐次线性约束（在 \(t=1\) 处取值为零 / 坐标和为零）削减了一个自由度，剩余一个二维子空间，由两个独立方向张成。对象从三维坐标换成多项式，结构丝毫不差。向量空间的语言让你只用一套推理同时分析 \(\R^n\)、多项式空间、函数空间——这恰恰是抽象的回报。

\subsection{诚实边界：向量空间不告诉你的事}

向量空间的公理框架非常干净，但它有意忽略了一些在实际问题中可能很重要的东西：

\begin{itemize}
\tightlist
\item \emph{长度和角度}：向量空间不要求定义``多长''或``多近''。两个多项式 ``\(t\)'' 和 ``\(t+0.0001\)'' 在向量空间结构中没有远近概念——引入内积之后才进入这个领域。那是正交性和最小二乘的地盘。
\item \emph{收敛和极限}：向量空间中无限个向量的线性组合（级数）不一定有定义。让分析意义上的极限在向量空间中成立，需要拓扑结构——这是泛函分析的工作。
\item \emph{乘法结构}：向量空间只管加法和数乘。矩阵之间有乘法，函数之间有点乘，但这些不是向量空间公理的一部分。需要乘法结构的场合（比如代数、环），会引入更丰富的代数系统。
\end{itemize}

这些不是向量空间的缺陷——正因为它刻意只管加法和数乘，它的结论才能适用于如此广泛的对象。更丰富的结构（内积、范数、乘法）是在向量空间公理之上逐层叠加的。学线性代数时不妨记住：现在你手里拿到的是一套在所有向量空间中都成立的``通用定理''。后面每多加一层结构，就多一组专用定理——线性代数这座大楼就是这样一层一层盖起来的。

延伸阅读：MIT 18.06：Column Space and Nullspace。
